\documentclass{article}
\usepackage[UTF8]{ctex}
\usepackage{amsmath,amssymb,amsfonts}
\usepackage{enumitem}
\usepackage{geometry}
\geometry{a4paper, margin=1in}
% === 错题本排版优化（不参与题目标识） ===
\setlength{\parindent}{0pt}
\setlength{\parskip}{0.35em plus 0.10em minus 0.08em}
\setlength{\emergencystretch}{3em}
\setlist{itemsep=0.25em, topsep=0.40em, parsep=0pt}
\allowdisplaybreaks[2]
% === 错题本排版优化结束 ===
\begin{document}
\section*{北京师范大学2024年高等代数真题}
\section*{1．北京师范大学2024年高等代数与空间解析几何试卷}
一．（10 分）设矩阵

\[
A=\left(\begin{array}{ccc}
2 & 0 & 0 \\
0 & 1 & -3 \\
0 & 2 & 1
\end{array}\right)
\]

求矩阵 \(B\) 使得 \(A+B=B A\) 。

二．（10 分）已知欧氏空间 \(\mathbb{R}^{3}\) 中一个基为 \(\alpha_{1}, \alpha_{2}, \alpha_{3}\) ，其度量矩阵为

\[
A=\left(\begin{array}{ccc}
1 & -1 & 0 \\
-1 & 2 & 0 \\
0 & 0 & 3
\end{array}\right)
\]

向量 \(\beta=\alpha_{1}+2 \alpha_{2}+3 \alpha_{3}\) ，求 \(\beta\) 长度．

三。（15 分）假设 \(A\) 是一个 \(n\) 阶反对称矩阵，其主对角线右上方元素全为 1 ，求 \(A\) 的秩和行列式。

四．（15 分）设

\[
\left\{\begin{array}{l}
x_{1}+x_{2}=a_{1} \\
x_{3}+x_{4}=a_{2} \\
x_{1}+x_{3}=b_{1} \\
x_{2}+x_{4}=b_{2}
\end{array}\right.
\]

是数域 \(F\) 上的线性方程组，试给出该方程组有解的充分必要条件，并在有解时求出其解．

五．（15 分）解答如下问题：\\

\subsection{题目1}
1．（7 分）设 \(f(x)\) 是实系数多项式，且 \(a+b \mathrm{i}\) 是 \(f(x)\) 的一个虚根，其中 \(a, b\) 是实数，证明 \(a-b \mathrm{i}\) 也是\(f(x)\) 的一个虚根．\\

\subsection{题目2}
2．（8 分）已知 \(\sqrt{2}+\mathrm{i}\) 是方程 \(x^{6}-3 x^{4}+11 x^{2}-9=0\) 的一个根，求该方程的其余根．

六．（20分）设矩阵

\[
A=\left(\begin{array}{ccc}
-2 & 0 & 0 \\
2 & a & 2 \\
3 & 1 & 1
\end{array}\right)
\]

且－1是 \(A\) 的一个特征值。\\

\subsection{题目3}
1．（4 分）求 \(a\) 的值。\\

\subsection{题目4}
2．（8 分）证明：\(A\) 可对角化，并求可逆矩阵 \(T\) 使得 \(T^{-1} A T\) 为对角阵。\\

\subsection{题目5}
3．（8 分）设 \(B\) 是一个3阶矩阵且满足 \(A B=B A\) ，证明：\(B\) 也可对角化．

七．（20 分）填空题．\\

\subsection{题目6}
1．（5 分）在 \(\triangle A B C\) 的三边 \(A B, B C, C A\) 上分别取三点 \(E, F, G\) 使得 \(A E: A B=\alpha, B F: B C=\beta\) ，\(C G: C A=\gamma\) ，那么 \(\triangle E F G\) 与 \(\triangle A B C\) 的面积之比为 \(\_\_\_\_\)。\\

\subsection{题目7}
2．（5 分）设曲线 \(a_{11} x^{2}+2 a_{12} x y+a_{22} y^{2}+2 b_{1} x+2 b_{2} y+c=0\) 为双曲线，则它的实轴和虚轴长度的平方和可以用二次曲线的不变量表达为 \(\_\_\_\_\)。\\

\subsection{题目8}
3．（5 分）双曲线 \(3 x^{2}-2 y^{2}+6 x+4 y=0\) 的渐近线为 \(\_\_\_\_\)。\\

\subsection{题目9}
4．（5分）过点（2，3，－4）且落在曲面 \(\frac{x^{2}}{4}+\frac{y^{2}}{9}-\frac{z^{2}}{16}=1\) 上的直线的标准方程为 \(\_\_\_\_\)．

八．（10 分）求过点 \((11,9,0)\) 与直线 \(\frac{x-1}{2}=\frac{y+3}{4}=\frac{z-5}{3}\) 和直线 \(\frac{x}{5}=\frac{y-2}{-1}=\frac{z+1}{2}\) 都相交的直线的标准方程．

九．（15 分）已知直线 \(l: \frac{x-1}{-3}=\frac{y+1}{1}=\frac{z-3}{-2}\) ，平面 \(\alpha: x-y+z-2=0\) ，直线 \(l\) 在平面 \(\alpha\) 上的投影直线记为 \(l_{1}\) ．\\

\subsection{题目10}
1．（7 分）试求直线 \(l_{1}\) 的方程．\\

\subsection{题目11}
2．（8 分）求直线 \(l\) 绕直线 \(l_{1}\) 旋转所得的曲面方程。

十．（20 分）把二次曲面方程 \(x^{2}+2 y^{2}+z^{2}-2 x y-2 y z+6 x+3 y+6 z+2=0\) 化简成标准方程，写出所用坐标变换，并说明这是哪种二次曲面。
\end{document}
